\begin{figure*}[t]
\centering
\scriptsize
\setlength{\fboxsep}{1em}
\setlength{\fboxrule}{1pt}
\fcolorbox{black}{white}{
\parbox{0.955\textwidth}{
\begin{alltt}


You will be provided instructions for completing a task followed by a task to complete.

Instructions:

1. Understand the question: Read the question carefully and understand what it's asking. The questions could be about a variety of topics in computer security, from technical aspects like buffer overflows and IP address spoofing, to more conceptual topics like message confidentiality and security principles.

2. Review the options: Look at each of the provided options. Even if you think you know the answer after reading the question, it's important to read all the options as there might be a better answer.

3. Apply your knowledge: Use your understanding of computer security to determine the correct answer. This might involve recalling specific facts or applying general principles.

4. Choose the best answer: Based on your understanding of the question and the options, choose the best answer. Remember, there might be more than one correct answer, but you need to choose the best one.

5. Justify your answer: Explain why you chose the answer you did. This will show that you didn't just guess, but actually used your knowledge of computer security to arrive at the answer.

6. Be confident: Even though this dataset is designed to be used in a zero-shot or few-shot setting, meaning you haven't seen any similar examples during training, you can still answer the questions correctly by applying your knowledge of computer security.

Question: What is penetration testing?

A. A procedure for testing libraries or other program components for vulnerabilities

B. Whole-system testing for security flaws and bugs

C. A security-minded form of unit testing that applies early in the development process

D. All of the above

Follow the instructions to generate an explanation that reasons towards the correct answer to the task above. End the explanation with the correct answer. [/INST] Explanation: 

1. Understand the question: The question is asking about the definition of penetration testing.

2. Review the options: Let's review the options provided:

A. A procedure for testing libraries or other program components for vulnerabilities - This option is incorrect as penetration testing is not limited to testing specific components, but rather the entire system.

B. Whole-system testing for security flaws and bugs - This option is correct as penetration testing involves testing the entire system for security vulnerabilities.

C. A security-minded form of unit testing that applies early in the development process - This option is incorrect as penetration testing is not a form of unit testing, and it is not limited to the development process.

D. All of the above - This option is incorrect as penetration testing is not limited to testing libraries or other program components, and it is not a form of unit testing.

3. Apply your knowledge: Based on our understanding of computer security, penetration testing is a process of testing a computer system, network, or web application to find vulnerabilities that an attacker could exploit. It involves simulating a cyber attack on the system to identify weaknesses and vulnerabilities.

4. Choose the best answer: Based on the options provided and our understanding of penetration testing, the best answer is B. Whole-system testing for security flaws and bugs.

5. Justify your answer: We chose option B because it accurately reflects the definition of penetration testing. Penetration testing involves testing the entire system for security vulnerabilities, which aligns with option B.

6. Be confident: We are confident in our answer because it is based on our understanding of computer security and the definition of penetration testing.

Answer: B. Whole-system testing for security flaws and bugs.

Therefore, the correct multiple choice label (just the letter) to the task is below.

Answer: B
\end{alltt}
}}
\caption{A complete running example of MMLU (Computer Security).}
\label{fig:mmlu_subject=computer_security_example}
\end{figure*}
